%\section{Puschkin -- Winterabend}
%\stepcounter{ctrpoeminyear}
%%\addcontentsline{toc}{section}{\thechapter.\arabic{ctrpoeminyear} \;}

\section*{Puschkin -- Winterabend}
%\clearpage
%\stepcounter{ctrpoeminyear}
\addcontentsline{toc}{section}{2026.Ü1 \;Puschkin -- Winterabend}
\fancyhead[LO,RE]{2026.Ü1}



\begin{strophe}
\vers{Der Sturm verhüllt den Himmel trübe,}
\vers{jagt wirbelnd Schnee durch kalten Wind;}
\vers{mal heult er auf gleich einem Tiere,}
\vers{mal weint er leise wie ein Kind,}
\vers{mal fährt er über unser altes,}
\vers{vom Wetter mürbes Strohdach hin,}
\vers{mal pocht er, wie hungrige Wand'rer,}
\vers{ans kleine Fenster zu uns drin.}
\end{strophe}

\begin{strophe}
\vers{Unser baufälliges Hüttchen}
\vers{steht traurig, dunkel, und so schwer}
\vers{Was bist auch du so still geworden,}
\vers{mein altes Mütterchen, so sehr?}
\vers{Hat dich des wilden Sturms Geheule,}
\vers{du treue Seele, müd' gemacht?}
\vers{Oder schläfst du zum zarten Säuseln}
\vers{deiner Spindel sacht umwacht?}
\end{strophe}

\begin{strophe}
\vers{Lass trinken auf den alten Kummer,}
\vers{oh, gute Freundnin, trink mit mir!}
\vers{Damit das Herze leichter schlummert --}
\vers{wo ist dein Becher? Bring ihn her!}
\vers{Sing mir das Lied der kleinen Meise,}
\vers{die selbst einst hinterm Meere sang.}
\vers{Sing mir das Lied der täglich Reise}
\vers{des jungen Mädchens Wassergang.}
\end{strophe}

\phantom{a}

\phantom{a}

\phantom{a}

\phantom{a}

\phantom{a}

\phantom{a}

\begin{strophe}
\vers{Der Sturm verhüllt den Himmel trübe,}
\vers{jagt wirbelnd Schnee durch kalten Wind;}
\vers{mal heult er auf gleich einem Tiere,}
\vers{mal weint er leise wie ein Kind.}
\vers{Lass trinken auf den alten Kummer,}
\vers{oh, gute Freundnin, trink mit mir!}
\vers{Damit das Herze leichter schlummert --}
\vers{wo ist dein Becher? Bring ihn her!}
\end{strophe}

\ttquotl Winterabend\ttquotr\ - Alexander Puschkin, 1825\\[0em]

Deutsche Übersetzung und lyrische Interpretation -

ChatGPT, 21.08.2026\\[0em]

Bearbeitung - Christoph Koloska, 21.08.2026

